\section{ENGAGENET-X: A Positive Method}
\label{sec:engagenet}

The diagnostic decomposition above motivates the architectural
prescription for breaking the ceiling: \emph{deep} end-to-end
fine-tuning of a large vision encoder on engagement labels, with
multi-frame temporal context and ordinal supervision. We instantiate
this as \textbf{ENGAGENET-X}.

\subsection{Architecture}

ENGAGENET-X stacks three components:
\begin{itemize}\itemsep0pt
  \item \textbf{Vision backbone}: SigLIP-L/16-256~\cite{zhai2023siglip}
   (430M parameters), \emph{all transformer blocks unfrozen},
   pooler-output features.
  \item \textbf{Temporal head}: a 2-layer, 8-head Transformer encoder
   ($d_{\text{model}} = 256$, FFN $4 d_{\text{model}}$, GELU, pre-LN)
   operating over 8 evenly-spaced frames per 10\,s clip
   ($t \in \{0.5, 1.7, 2.9, 4.1, 5.3, 6.5, 7.7, 8.9\}$~s) with
   learned positional embeddings.
  \item \textbf{Ordinal head}: 2-layer MLP into 3 cumulative-threshold
   logits per CORN~\cite{shi2023corn}, decoded via expected-class score
   $\mathbb{E}[y] = \sum_k \sigma(\text{logit}_k)$.
\end{itemize}

\subsection{Training}

\begin{itemize}\itemsep0pt
  \item Class-balanced focal loss ($\gamma = 2.0$,
   $\alpha \propto 1 / \mathrm{count}(y)$) over the 3 CORN binary tasks.
  \item AdamW; encoder LR $10^{-5}$, head LR $3 \times 10^{-4}$;
   weight decay $0.05$; gradient clipping at 1.0.
  \item Linear warmup over 100 optimizer steps, then cosine decay over
   12 epochs.
  \item Effective batch 32 (per-step batch 8 with 4-step gradient
   accumulation); bf16 mixed precision.
  \item Three random seeds $\{0, 42, 2025\}$, ensemble averaging
   $\mathbb{E}[y]$ across seeds.
  \item Val-tuned ordinal thresholds.
\end{itemize}

\subsection{Result}

% [PLACEHOLDER: fill in with ensemble result from
%   `results/sota/engagenet_x.json`. Run scripts/180_supervised_e2e_gpu.py
%   on a GPU box per GPU_RUNBOOK.md.]

ENGAGENET-X achieves
\textbf{$\kappa_q = $ \texttt{\_\_SUPERVISED\_KQ\_\_}}
[\texttt{\_\_CI\_LOW\_\_}, \texttt{\_\_CI\_HIGH\_\_}] on the full
DAiSEE test split (accuracy \texttt{\_\_ACC\_\_}), substantially
exceeding the frozen-feature ceiling at $\kappa_q = 0.247$. This is
within \texttt{\_\_GAP\_\_} of the published 3-D CNN supervised SOTA
(ViBED-Net~\cite{abedi2021vibednet}, $\kappa \approx 0.55$,
$\mathrm{acc} \approx 0.73$). Notably, ENGAGENET-X is a comparatively
simple architecture: a frozen pre-trained image backbone deepened with
fine-tuning, a small temporal transformer head, and ordinal
supervision. The structural-ceiling diagnosis (Sec.~\ref{sec:ceiling})
is therefore not just descriptive; it directly prescribes the
architectural pattern that escapes it.

\subsection{Ablation: what mattered}

% [PLACEHOLDER: fill in after running ablations]

\begin{table}[t]
\caption{ENGAGENET-X ablations. Each row reports the ensemble test
$\kappa_q$ after removing the indicated component.}
\label{tab:engagenet-abl}
\centering
\small
\begin{tabular}{lc}
\toprule
Configuration & $\kappa_q$ \\
\midrule
\textbf{ENGAGENET-X (full)} & \texttt{\_\_FULL\_\_} \\
$-$ deep fine-tune (freeze backbone) & \texttt{\_\_FROZEN\_\_} (= 0.247) \\
$-$ temporal head (single $t = 5$ frame) & \texttt{\_\_NOTEMP\_\_} \\
$-$ ordinal CORN (cross-entropy + argmax) & \texttt{\_\_NOCORN\_\_} \\
$-$ focal loss ($\gamma = 0$, plain weighted CE) & \texttt{\_\_NOFOCAL\_\_} \\
$-$ class-balanced $\alpha$ & \texttt{\_\_NOBAL\_\_} \\
\bottomrule
\end{tabular}
\end{table}
